\section{Introduzione}

Questo report documenta il lavoro svolto nel progetto del corso ISW2, articolato in quattro
milestone progressive sull'analisi della relazione tra code smell, bugginess e predizione
dei difetti nel progetto open source \textbf{Apache OpenJPA} (framework JPA della Apache
Software Foundation). I dati sui ticket di bug sono stati recuperati tramite le API REST
di Jira (\texttt{issues.apache.org}); il codice sorgente è stato estratto dal repository
Git ufficiale.

\begin{enumerate}
    \item \textbf{Milestone 1 -- Costruzione del dataset}: raccolta e sanitizzazione dei
    ticket Jira, estrazione degli snapshot di classi Java per release, calcolo di 19
    metriche di processo e prodotto, etichettatura della bugginess.

    \item \textbf{Milestone 2 -- Identificazione di modelli accurati}: feature selection,
    data balancing, addestramento e valutazione di classificatori per la predizione della
    bugginess tramite 10-times 10-fold cross-validation (27--60 combinazioni).

    \item \textbf{Milestone 3 -- Analisi dell'impatto dei code smell}: analisi controfattuale
    per quantificare quante classi buggy sarebbero state prevenute in assenza di code smell,
    con studio di ridondanza (correlazione Spearman e ablation study).

    \item \textbf{Milestone 4 -- Refactoring automatico}: valutazione della possibilità
    di eliminare automaticamente i code smell tramite refactoring guidato, sotto diverse
    condizioni di disponibilità di test.
\end{enumerate}
